\documentclass[10pt,a4paper]{amsart}
\usepackage[margin=19mm]{geometry}
\usepackage{amsmath,amssymb,mathtools}
\usepackage[scr=rsfs,cal=euler]{mathalfa}
\usepackage{xfrac}
\usepackage[unicode,hidelinks,bookmarks=false]{hyperref}
\newcommand{\ie}{i.e.\ }
\newcommand{\cf}{cf.\ }
\newcommand{\resp}{respectively\ }
\newcommand{\Subsection}{\subsection}
\providecommand{\nobreakdash}{\nobreak\hbox{-}}
\setlength{\emergencystretch}{2em}
\multlinegap0pt
\usepackage{bm}
\usepackage{bbm}
\usepackage{dsfont}
\usepackage{xspace}
\usepackage{cancel}
\usepackage{mathtools}
\usepackage{amsthm}
\makeatletter
\@ifundefined{theorem}{\newtheorem{theorem}{Theorem}}{}
\@ifundefined{claim}{\newtheorem{claim}[theorem]{Claim}}{}
\@ifundefined{lemma}{\newtheorem{lemma}[theorem]{Lemma}}{}
\@ifundefined{obs}{\newtheorem{obs}[theorem]{Observation}}{}
\makeatother
\title[An Upper Bound on the Hat Guessing Number of Graphs]{An Upper Bound on the Hat Guessing Number of Graphs}
\author{Mason Shurman, Scott Albert Sibley}
\date{Source: arXiv:2607.07994v1 (CC BY 4.0)}
\begin{document}
\maketitle

\noindent\emph{Source and licence.} An Upper Bound on the Hat Guessing Number of Graphs. Version arXiv:2607.07994v1 (CC BY 4.0), distributed under the Creative Commons Attribution 4.0 International licence, as recorded in the arXiv licence metadata for this exact version and shown on the abstract page https://arxiv.org/abs/2607.07994v1. Full rights notice: licenses/arxiv-2607.07994-CC-BY-4.0.txt.\par\smallskip

\noindent\textit{Editorial scope.} Counted unit: the Lemma whose source label is \texttt{3.2lem} in the article (source0.tex lines 316--318, source label \texttt{3.2lem}), delivered together with the article's own complete original proof of it (source0.tex lines 319--346, one \texttt{proof} environment of 28 lines). No mathematics is written, repaired or re-derived: statement and proof are the article's own text line for line. Required context retained uncounted: 3.1obs:degs (lines 314--314). Every definition, display and label of those lines is retained unchanged. Omitted from this package and not redistributed: 1--313, 315--315, 347--526. Nothing inside a retained excerpt is omitted or altered: every retained body is delivered exactly as the article's lines are written, so no omission receipt of any kind accompanies this package. Citations: the retained excerpts carry no citation command at all, so no bibliography entry of the article is transcribed and no reference is redistributed. Reference adaptation: none was needed and none was made; every reference inside the delivered unit resolves inside it, so no unresolved reference or citation is produced. Printed slips kept literal: no printed word, symbol, hypothesis or number of the counted statement or of its proof was altered. Layout and class: the article's own class is \texttt{amsart} with packages including graphicx, inputenc, amsfonts, amsmath, amssymb, amsthm, dsfont, geometry, comment, pgffor, pgfplots, mathtools, enumerate, lineno; the wrapper is the lane's pinned \texttt{amsart} editorial wrapper at 10pt on A4, which restates the theorem-like environments the retained text needs and adds the article's own macro definitions that the retained text needs (none), with extra wrapper packages (bm, bbm, dsfont, xspace, cancel, mathtools). The delivered numbering is the unit's own. Assets: no retained span contains a figure environment, a graphics-inclusion command, a PDF-page inclusion, a table or a TikZ picture; no image file is copied into this package, the package declares no asset dependency and no image file is redistributed with it. Licence: version arXiv:2607.07994v1 of this article is distributed under the Creative Commons Attribution 4.0 International licence, as recorded in the arXiv licence metadata for this exact version and shown on the abstract page https://arxiv.org/abs/2607.07994v1. A full rights notice is delivered at licenses/arxiv-2607.07994-CC-BY-4.0.txt. Characters: every retained excerpt is pure ASCII, so the delivered engine, XeLaTeX with its default Latin Modern text font, reports no missing character.
\begin{obs}\label{3.1obs:degs} $$e(H)=\sum_{c\in X_1} \deg(c)=\sum_{c' \notin X_1} \deg(c').$$\end{obs}

\begin{lemma}\label{3.2lem} %(Double-counting bound)
     $$(n-1-\Delta)|X_1| \leq e(H) \leq \sum_{i=2}^n i(k-1)|X_i|.$$
\end{lemma}

\begin{proof}
%To prove this, we prove each part separately, using the definition of $H$. 
Let $c \in X_1$, $c' \not \in X_1$. We will prove the lower and upper bounds by using Observation \ref{3.1obs:degs} and bounding $\deg (c)$ and $\deg(c')$ respectively.
\begin{claim}\label{3.3clm} For any $c\in X_1$, 
    $$\deg (c)\geq n-1-\Delta.$$
\end{claim}
\begin{proof}
    Let $c \in X_1$, and let $v \in V(G)$ be the vertex that guesses its color correctly in $c$. Let $u \in V(G)\setminus \{v\}$ be a vertex not adjacent to $v$, and let $a_u$ be the color that $u$ guesses under the coloring $c$. 
        Let $c_u \in \mathcal{C}$ be the coloring where:
        $$c_u(x):= \begin{cases} a_u \text{ if }x = u \\c(x) \text{ otherwise.}
        \end{cases}$$
        Then $c$ and $c_u$ are adjacent in $H$. For each $u \in V(G)$ not adjacent to $v$, the colorings $c_u$ are distinct, and there are at least $n-1-\Delta$ vertices not adjacent to $v$, not including $v$ itself. Thus, $\deg(c)\geq n-1-\Delta$. %\footnote{Note that this describes all the edges of any $c\in X_1$, so replacing $\Delta$ with $|N(v)|$ gives an equality. This may be useful for graphs of small average degree, but large maximum degree.} 
\end{proof}
\begin{claim}\label{3.4clm} For any $c'\in X_i$, $i\neq 1$,  
    $$\deg(c')\leq  i(k-1).$$
\end{claim}
\begin{proof}
    Let $c' \in X_i$, $i\neq 1$, and let $c \in X_1$ be adjacent to $c'$. Then there exist nonadjacent vertices $u$ and $v$ such that $v$ guesses correctly in $c$ (and $c'$), $u$ guesses correctly in $c'$, and $c'(w)=c(w)$ for all vertices $w \neq u$.
    In fact, an edge into $c'$ is determined by this unique vertex $u$ where the coloring changes, and that vertex must guess correctly in $c'$. 

    
    For any fixed $c' \in X_i$, there are $i$ vertices that guess correctly in $c'$, so $i$ potential choices for $u$. For each of these vertices, there are $k-1$ ways to change its color. Since any $c$ adjacent to $c'$ must agree with $c'$ on every vertex except $u$, there are at most $i(k-1)$ colorings $c$ adjacent to $c'$. %\mason{Can we improve this writing?}
\end{proof}

Putting Claims \ref{3.3clm} and \ref{3.4clm} together with Observation \ref{3.1obs:degs}, we get % \ref{3.2lem}, since 
$$e(H)=\sum_{c\in X_1}  \deg (c)\geq (n-1-\Delta)|X_1|,$$ and $$e(H)=\sum_{c' \notin X_1} \deg(c')\leq \sum _{c' \notin X_1, c' \in X_i} i(k-1)=\sum i(k-1)|X_i|.$$
This completes the proof of Lemma \ref{3.2lem}.
\end{proof}

\end{document}
